\documentclass[10pt,a4paper]{amsart}
\usepackage[margin=19mm]{geometry}
\usepackage{amsmath,amssymb,mathtools}
\usepackage[scr=rsfs,cal=euler]{mathalfa}
\usepackage{xfrac}
\usepackage[unicode,hidelinks,bookmarks=false]{hyperref}
\newcommand{\ie}{i.e.\ }
\newcommand{\cf}{cf.\ }
\newcommand{\resp}{respectively\ }
\newcommand{\Subsection}{\subsection}
\providecommand{\nobreakdash}{\nobreak\hbox{-}}
\setlength{\emergencystretch}{2em}
\multlinegap0pt
\usepackage{bm}
\usepackage{bbm}
\usepackage{dsfont}
\usepackage{xspace}
\usepackage{cancel}
\usepackage{mathtools}
\usepackage{amsthm}
\makeatletter
\@ifundefined{theorem}{\newtheorem{theorem}{Theorem}}{}
\@ifundefined{lemma}{\newtheorem{lemma}[theorem]{Lemma}}{}
\makeatother
\makeatletter
\expandafter\ifx\csname absolutelynopagebreak\endcsname\relax
\newenvironment{absolutelynopagebreak}
  {\par\nobreak\vfil\penalty0\vfilneg
   \vtop\bgroup}
  {\par\xdef\tpd{\the\prevdepth}\egroup
   \prevdepth=\tpd}
\fi
\expandafter\ifx\csname clproof\endcsname\relax
\newenvironment{clproof}{\par
	\pushQED{\cqed}%
	\normalfont \topsep6\p@\@plus6\p@\relax
	\trivlist
	\item\relax
	{\itshape
		Proof of the claim\@addpunct{.}}\hspace\labelsep\ignorespaces
}{%
	\hfill\popQED\endtrivlist\@endpefalse
}
\fi
\makeatother
\title[On the generalized coloring numbers]{On the generalized coloring numbers}
\author{Sebastian Siebertz}
\date{Source: arXiv:2501.08698v3 (CC BY 4.0)}
\begin{document}
\maketitle

\noindent\emph{Source and licence.} On the generalized coloring numbers. Version arXiv:2501.08698v3 (CC BY 4.0), distributed under the Creative Commons Attribution 4.0 International licence, as recorded in the arXiv licence metadata for this exact version and shown on the abstract page https://arxiv.org/abs/2501.08698v3. Full rights notice: licenses/arxiv-2501.08698-CC-BY-4.0.txt.\par\smallskip

\noindent\textit{Editorial scope.} Counted unit: the Lemma whose source label is \texttt{thm:augmentationneighborhood} in the article (source0.tex lines 2179--2207, source label \texttt{thm:augmentationneighborhood}), delivered together with the article's own complete original proof of it (source0.tex lines 2208--2255, one \texttt{proof} environment of 48 lines). No mathematics is written, repaired or re-derived: statement and proof are the article's own text line for line. Required context retained uncounted: none. Every definition, display and label of those lines is retained unchanged. Omitted from this package and not redistributed: 1--2178, 2256--4358. Nothing inside a retained excerpt is omitted or altered: every retained body is delivered exactly as the article's lines are written, so no omission receipt of any kind accompanies this package. Citations: the retained excerpts carry no citation command at all, so no bibliography entry of the article is transcribed and no reference is redistributed. Reference adaptation: none was needed and none was made; every reference inside the delivered unit resolves inside it, so no unresolved reference or citation is produced. Printed slips kept literal: no printed word, symbol, hypothesis or number of the counted statement or of its proof was altered. Layout and class: the article's own class is \texttt{elsarticle} with packages including anyfontsize, stix, lineno, amsfonts, amsmath, amsthm, mathrsfs, enumerate, graphics, mathtools, microtype, nowidow, xcolor, babel; the wrapper is the lane's pinned \texttt{amsart} editorial wrapper at 10pt on A4, which restates the theorem-like environments the retained text needs and adds the article's own macro definitions that the retained text needs (none), with extra wrapper packages (bm, bbm, dsfont, xspace, cancel, mathtools). The delivered numbering is the unit's own. Assets: no retained span contains a figure environment, a graphics-inclusion command, a PDF-page inclusion, a table or a TikZ picture; no image file is copied into this package, the package declares no asset dependency and no image file is redistributed with it. Licence: version arXiv:2501.08698v3 of this article is distributed under the Creative Commons Attribution 4.0 International licence, as recorded in the arXiv licence metadata for this exact version and shown on the abstract page https://arxiv.org/abs/2501.08698v3. A full rights notice is delivered at licenses/arxiv-2501.08698-CC-BY-4.0.txt. Characters: every retained excerpt is pure ASCII, so the delivered engine, XeLaTeX with its default Latin Modern text font, reports no missing character.
\begin{lemma}\label{thm:augmentationneighborhood}
  Let~$G$ be a graph, let~$r$ be a positive integer and 
  let $w$ be an $r$-fraternity function of~$G$.  
  Let~$u,v\in V(G)$
  and assume that \mbox{$u=v_0, v_1,\ldots, v_\ell=v$} is a path of
  length $\ell\leq r$ between~$u$ and~$v$ in~$G$. Then 
  one of the following holds. 
  \begin{enumerate}
	\item There are indices $0< i_1<\ldots<i_k< \ell$ 
	such that $u,v_{i_1}, \ldots, v_{i_k},v$ is a directed 
	path from $u$ to $v$ in $G^w_{\leq r}$ and for all 
  $0\leq i_m<i_n\leq \ell$ we have $w(v_{i_m}, v_{i_n})\leq 
  i_n-i_m$.
    \item There are indices $0< i_1<\ldots<i_k< \ell$ 
	such that $v,v_{i_k}, \ldots, v_{i_1},u$ is a directed 
	path from $v$ to $u$ in $G^w_{\leq r}$ and for all 
  $0\leq i_m<i_n\leq \ell$ we have $w(v_{i_n}, v_{i_m})\leq 
  i_n-i_m$.
    \item There are indices $0< i_1<\ldots<i_k< \ell$ and 
    an index $i_j\in \{i_1,\ldots, i_k\}$ such that 
    $u, v_{i_1},\ldots, v_{i_j}$ and $v, v_{i_k},\ldots, v_{i_j}$
    are directed paths in $G^w_{\leq r}$ from $u$ to $v_{i_j}$ and from 
    $v$ to $v_{i_j}$, respectively. Furthermore, for all 
  $0\leq i_m<i_n\leq i_j$ we have $w(v_{i_m}, v_{i_n})\leq 
  i_n-i_m$ and for all 
  $i_j\leq i_m<i_n\leq \ell$ we have $w(v_{i_n}, v_{i_m})\leq 
  i_n-i_m$.
  \end{enumerate}
\end{lemma}

\begin{proof}
We prove the statement by induction on $\ell$. For $\ell=1$, 
we have $(u,v)\in E(\vec{G}^w_{\leq 1})$ if and only if 
$w(u,v)=1$, hence, one of the first two items is true. 
Now assume that 
$\ell>1$ and assume the statements hold for 
$v_0,v_1,\ldots, v_{\ell-1}$. Assume we are in the first case and 
there are indices $0< i_1<\ldots<i_k< \ell-1$ 
such that $u,v_{i_1}, \ldots, v_{i_k},v_{\ell-1}$ is a directed 
path from $u$ to $v_{\ell-1}$ in $G^w_{\leq r}$ and that 
for all $0\leq i_m<i_n\leq \ell-1$ we have 
$w(v_{i_m}, v_{i_n})\leq i_n-i_m$. 
We distinguish
two cases. If $w(v_{\ell-1},v)=1$, $u,v_{i_1}, \ldots, 
v_{i_k},v_{\ell-1},v$ is a directed path in $G^w_r$, hence the
first item holds for $u=v_0,v_1,\ldots, v_\ell=v$. If $w(v,v_{\ell-1})=1$, then we set~$v_j$ in 
the third item to $v_{\ell-1}$ and the third item holds for 
$u=v_0,v_1,\ldots, v_\ell=v$. The statements about the 
weight function follow directly from the hypothesis. 

Now assume we are in the second case and there are indices 
$0< i_1<\ldots<i_k< \ell-1$ 	such that $v_{\ell-1},v_{i_k}, 
\ldots, v_{i_1},u$ is a directed path from $v_{\ell-1}$ to $u$ 
in $G^w_{\leq r}$ and that for all $0\leq i_m<i_n\leq \ell-1$ we have 
$w(v_{i_n},v_{i_m})\leq i_n-
i_m$. Again we distinguish two cases. If $w(v,v_{\ell-1})=1$, then 
the first case holds for \mbox{$u=v_0,v_1,\ldots, v_\ell=v$}. The 
additional statement about the weight function follows directly
from the hypothesis. If $w(v_{\ell-1},v)=1$, then we have 
$w(v,v_{i_k})\leq \ell-i_k$ or $w(v_{i_k},v)\leq \ell-i_k$, 
as~$w$ is a fraternity function and by assumption, we have 
$w(v_{\ell-1},v_{i_k})\leq \ell-i_k-1$. In the former case 
we have a directed path $v,v_{i_k}, 
\ldots, v_{i_1},u$ from $v$ to $u$ in $G^w_{\leq r}$, that is, 
the second case holds for $u=v_0,v_1,\ldots, v_\ell=v$. 
In the latter case, we apply the argument again to the sequence 
$v,v_{i_k}, \ldots, v_{i_1},u$ ($v_{i_k}$ taking the role of 
$v_{\ell-1}$). Again, we either get a directed path $v,v_{i_{k-1}}, 
\ldots, v_{i_1},u$ from $v$ to $u$ in $G^w_{\leq r}$, or
we apply the argument again. In some step, we either get
a directed path $v,v_{i_{m}}, \ldots, v_{i_1},u$ from 
$v$ to $u$ in $G^w_{\leq r}$ for some $1\leq i\leq k$, or 
a directed arc $(u,v)$ in~$G^w_{\leq r}$. The additional property on 
the weight function is easily verified in each step. 

The argument in the third case is analogous to the argument 
in the second case. 
\end{proof}

\end{document}
